\section{Cục bộ rẻ hơn toàn cục, và cái giá của phép so}
\label{sec:ch13-cuc-bo-va-toan-cuc}

\index{lời giải thích cục bộ!chi phí mô tả|(}%
Mọi chặn của mục trước đều tính cho lời giải thích toàn cục, tức một hàm $g$
phải bám $f$ trên cả không gian đầu vào. Nhưng mọi phương pháp mà Phần~II dựng
đều là phương pháp cục bộ: LIME khớp một mô hình thay thế trong lân cận một điểm
và SHAP chia công cho một dự đoán đơn lẻ. Mục này đọc phần bài báo dành cho chỗ
khác nhau ấy, và đó cũng là chỗ duy nhất trong 65 trang mà bài nhắc tới hai
phương pháp trên.

Bài định nghĩa sai số cục bộ đúng như người ta chờ đợi: cho một điểm $x_0$, một
bán kính $r$ và một ngân sách, lấy chỗ lệch lớn nhất giữa $f$ và $g$ trong quả
cầu bán kính $r$ quanh $x_0$, rồi lấy cận dưới đúng trên mọi $g$ trong ngân
sách~\autocite{p26limits}. Từ đó ra ngay một bổ đề mà bài ghi cho đủ: lấy trung
bình sai số cục bộ ấy trên mọi điểm $x_0$ thì được một đại lượng không vượt quá
sai số toàn cục, vì cận trên trong một quả cầu không lớn hơn cận trên trên cả
không gian~\autocite{p26limits}. Giải thích cục bộ không khó hơn giải thích toàn
cục, và câu ấy đúng vì lý do hình thức chứ không vì lý do nào sâu hơn.

Kết quả thật nằm ở định lý tiếp theo, và phát biểu của nó có một mệnh đề mà mục
này xoay quanh. Với $f$ là hàm $L$-Lipschitz và $N(x_0)$ là một lân cận nằm gọn
trong quả cầu bán kính $r$, bài chứng minh rằng độ phức tạp cục bộ bằng
$\mathcal{O}(1)$ khi $\delta \ge Lr$ và bằng $\mathcal{O}(d \log(Lr/\delta))$
khi $\delta < Lr$; và bài viết rõ đại lượng ấy là độ dài mô tả ngắn nhất của một
chương trình tính được lời giải thích với sai số dưới $\delta$ trong lân cận ấy,
\emph{khi chương trình được truy cập $f$ như một tiên
tri}~\autocite{p26limits}.

Trường hợp thứ nhất dễ thấy và đáng viết ra vì nó cho biết cả hai vế đến từ đâu.
Nếu ngưỡng sai số $\delta$ đã lớn hơn mức biến thiên tối đa $Lr$ của $f$ trong
quả cầu, thì hàm hằng $g(x) = f(x_0)$ đã đủ, vì tính Lipschitz cho
$|f(x) - f(x_0)| \le Lr \le \delta$ với mọi $x$ trong quả cầu. Chương trình mô
tả $g$ chỉ cần nói: trả về $f(x_0)$. Nó không cần chứa giá trị $f(x_0)$, vì tiên
tri cung cấp giá trị ấy khi được hỏi. Đó là lý do độ dài mô tả là hằng số.
Trường hợp thứ hai làm cùng một việc với một lưới nhiều ô thay vì một ô, và các
giá trị $f$ tại các ô cũng lấy từ tiên tri.

Cuối phần chứng minh, bài so hai đại lượng và rút ra câu mà chương này phải đọc
kỹ. Bài viết rằng với lời giải thích toàn cục trên một miền đường kính $D$, độ
phức tạp có dáng $\mathcal{O}\big((LD/\delta)^d \, d \log(LD/\delta)\big)$, tức
mũ theo $d$, trong khi độ phức tạp cục bộ chỉ là hàm lôgarit của tỉ số
$Lr/\delta$, nên lời giải thích cục bộ có thể đơn giản hơn lời giải thích toàn
cục theo hàm mũ~\autocite{p26limits}. Rồi ngay dưới đó, cách một nhận xét, bài
kết luận: điều đó lý giải vì sao các phương pháp giải thích cục bộ như LIME và
SHAP dễ xử lý hơn các phương pháp toàn cục~\autocite{p26limits}.

Nhận xét của sách, không phải của bài: phép so ấy đặt cạnh nhau hai đại lượng
không đo cùng một cách. Vế cục bộ là độ dài mô tả tương đối, đo khi đã có tiên
tri; vế toàn cục là độ dài mô tả tuyệt đối, không có tiên tri nào. Mà số hạng
trội trong mọi chặn toàn cục của mục~\ref{sec:ch13-tran-cao-bao-nhieu} chính là
chi phí mã hóa các giá trị hàm tại các ô lưới, tức đúng phần mà vế cục bộ được
tiên tri cho không. Hiệu số mũ giữa hai vế vì thế phần lớn là hiệu giữa hai quy
ước đo, chứ không phải giữa hai bài toán. Bài tự phát biểu quy ước ấy hai trang
trước đó, trong chính phát biểu của định lý, và không mang nó theo vào phép so.

Đây không phải một mệnh đề sai và sách không gọi nó là sai. Định lý về độ phức
tạp cục bộ đúng như đã phát biểu, với quy ước như đã phát biểu; chỗ hụt nằm ở
câu so sánh và ở câu kết luận rút từ câu so sánh ấy. Còn một chỗ hụt nhỏ hơn
trong cùng phần chứng minh, ghi lại cho đủ: sau khi đếm được
$\mathcal{O}\big((Lr/\delta)^d \, d\log(Lr/\delta)\big)$ bit cho phép phủ, bài
nói có thể cải thiện xuống $\mathcal{O}(d \log(Lr/\delta))$ bằng cách chia đôi
đệ quy, với lý do là độ sâu đệ quy cỡ $\log(Lr/\delta)$ và mỗi mức tốn
$\mathcal{O}(d)$ bit~\autocite{p26limits}. Phép đếm ấy là chi phí mô tả một
đường đi xuống cây chia đôi, không phải chi phí mô tả cả cách phân hoạch, và
phát biểu của định lý lấy con số sau khi cải thiện.

\index{lời giải thích cục bộ!chi phí mô tả và độ trung thực}%
Nhưng giả sử cứ nhận phép so ấy nguyên vẹn, thì nó nói gì với Phần~IV? Nó nói
rằng khớp một hàm đơn giản với $f$ trong một lân cận nhỏ là việc rẻ. Điều ấy
đúng, và nó đúng vì một lý do mà chương~\ref{ch:lime} đã dựng sẵn:
mục~\ref{sec:ch03-khop-mo-hinh-thay-the} cho thấy LIME tối ưu hóa đúng đại lượng
đó theo thiết kế, bằng cách cực tiểu một hàm mất mát có trọng số theo lân cận.
Một định lý nói rằng đại lượng ấy nhỏ được thì xác nhận thiết kế của LIME chứ
không xác nhận thêm điều gì. Còn câu hỏi của định
nghĩa~\ref{def:ch01-do-trung-thuc}, tức lời giải thích ấy có phản ánh đúng cách
mô hình đi tới đầu ra hay không, thì không nằm trong phát biểu nào của mục này.
Chương~\ref{ch:lime-nao-dang-tin} đã đọc bốn mươi tám biến thể của LIME, hai
phần ba trong số đó nhắm vào độ khớp hoặc \tn{tính ổn định}{stability}, và ghi
rằng lĩnh vực
không có cách xếp hạng chúng; một cái trần đặt lên độ khớp không đụng tới chỗ
đó.

\index{lời giải thích cục bộ!chi phí mô tả|)}%
